% Generated draft: requires mathematical and source review.
\documentclass[11pt]{article}
\usepackage[margin=25mm]{geometry}
\usepackage{fontspec}
\setmainfont{Latin Modern Roman}
\usepackage{amsmath,amssymb,mathtools,unicode-math}
\setmathfont{Latin Modern Math}
\usepackage{enumitem,needspace,xurl,hyperref}
\setlength{\parindent}{0pt}
\setlength{\parskip}{5pt}
\setlength{\emergencystretch}{3em}
\newcommand{\N}{\mathbb{N}}
\newcommand{\Z}{\mathbb{Z}}
\newcommand{\Q}{\mathbb{Q}}
\newcommand{\R}{\mathbb{R}}
\newcommand{\C}{\mathbb{C}}
\newcommand{\F}{\mathbb{F}}
\newcommand{\A}{\mathbb{A}}
\newcommand{\PP}{\mathbb{P}}
\newcommand{\E}{\mathbb{E}}
\newcommand{\eps}{\varepsilon}
\newcommand{\dd}{\,\mathrm{d}}
\begin{document}
\section*{UnsolvedMath ID: 30002521}
\textbf{Source problem code:} catalogue entry
\subsection{Definitions and mathematical statement}
The supplied record does not specify the base field or the precise convention for periodicity. The following is a provisional formulation using ordinary syzygy periodicity; these conventions and any additional source hypotheses require confirmation.

Let \(k\) be a field and let \(A\) be a finite-dimensional associative unital \(k\)-algebra. Write \(\operatorname{mod}A\) for the category of finite-dimensional right \(A\)-modules, choosing right modules as a notation convention. For a module \(M\), let \(\Omega_A(M)\) be the kernel of a projective cover, and define higher syzygies by iteration. Call a nonzero module \(M\) periodic if
\[
\exists n\in\Z_{>0}\qquad \Omega_A^n(M)\cong M
\]
as \(A\)-modules. The period may depend on \(M\).

Put \(A^{\mathrm e}=A\otimes_k A^{\mathrm{op}}\), where \(A^{\mathrm{op}}\) is the opposite algebra. Regard the regular bimodule \(A\) as a left \(A^{\mathrm e}\)-module through \((a\otimes b^{\mathrm{op}})x=axb\). Call \(A\) a periodic algebra if
\[
\exists p\in\Z_{>0}\qquad \Omega_{A^{\mathrm e}}^p(A)\cong A
\]
as bimodules, with syzygies taken in a minimal projective bimodule resolution.

\begin{enumerate}
\item Determine whether, for the intended class of finite-dimensional algebras, the implication
\[
\left(\forall\text{ simple }S\in\operatorname{mod}A\;\exists n_S\in\Z_{>0}:\ \Omega_A^{n_S}(S)\cong S\right)
\ \Longrightarrow\ 
\left(\exists p\in\Z_{>0}:\ \Omega_{A^{\mathrm e}}^p(A)\cong A\right)
\]
holds. No common period for the simple modules is assumed in the supplied statement.
\item Classify periodic finite-dimensional algebras up to Morita equivalence, and separately up to derived equivalence. Under the usual fixed-field interpretation, Morita equivalence means an equivalence of module categories as \(k\)-linear categories, and derived equivalence means a \(k\)-linear triangulated equivalence
\[
D^b(\operatorname{mod}A)\simeq D^b(\operatorname{mod}B),
\]
where \(D^b\) denotes the bounded derived category.
\end{enumerate}

The record does not determine whether \(k\) must be algebraically closed or satisfy other restrictions, whether further hypotheses on \(A\) are intended, or whether periodicity is instead meant in a stable category. This distinction affects boundary cases: under the displayed ordinary-syzygy definition a nonzero projective module is not periodic, whereas it represents the zero object in the stable category. Consequently, the treatment of projective simple modules and semisimple algebras requires confirmation. The fixed-field and \(k\)-linear conventions for the classification questions also require confirmation.
\subsection{Short English statement}
Ask whether periodicity of every simple module forces a finite-dimensional algebra to be periodic as a bimodule over itself, and classify periodic algebras up to Morita equivalence and, separately, derived equivalence. The base-field hypotheses and periodicity conventions are unresolved in the supplied record, so the displayed formulation is provisional.
\subsection{Sources}
\begin{itemize}
\item Local registry: \texttt{data/unsolvedmath/problems.json}, supplied record with permanent UnsolvedMath ID \(30002521\), title \emph{Periodic Simple Modules and Periodic Algebras}, source problem code \texttt{catalogue entry}. The supplied collection is \emph{Oberwolfach Reports Open Problems}; the proposed year is 2014, and no individual proposer is identified.
\item Supplied workshop provenance: \emph{Representation Theory of Quivers and Finite Dimensional Algebras} (2014), Oberwolfach Reports, pp.\ 453--529, DOI \texttt{10.4171/owr/2014/08}. Also listed in the literature sources as \emph{Oberwolfach Report 08/2014}: \url{https://doi.org/10.4171/owr/2014/08}. The supplied background describes Oberwolfach Reports as documenting mathematical workshops and their open-problem discussions.
\item Supplied literature source: \emph{Periodic resolutions and self-injective algebras}, \url{https://doi.org/10.1016/j.jpaa.2009.09.012}.
\item Supplied literature source: \emph{Periodic trivial extension algebras}, \url{https://arxiv.org/abs/2012.11927}.
\item Supplied literature assessment, checked \texttt{2026-08-21}: the record labels the problem open and reports that both the implication and the requested classification remain open in general, despite proofs for important families such as group blocks and finite representation type. This is supplied dated curation derived from the OWR material, not an independently verified assessment. The scope and hypotheses of the claimed special-family results are not supplied; time-sensitive status and cited claims require independent verification.
\item The supplied statement status is \texttt{corrected\_verified}, accompanied by a claim of manual semantic, typographic, and LaTeX review of a meaning-preserving cleanup. This is a supplied verification label, not independent mathematical verification of this draft or its conventions.
\end{itemize}
\end{document}
